\chapter{Aggregate Supply and the Phillips Curve}\label{ch:b1:aggregate-supply}

In the United States from 1960 to 2018, the years of low unemployment were
years in which inflation rose, and the years of high unemployment years in
which it fell. In 1974, with unemployment at $5.6\%$, inflation rose by $4.82$
points; in 1982, at $9.7\%$, it fell by $4.16$. % figdata: accel_points
(Inflation here is the percent change of the yearly average consumer price
index.)
The level of inflation in a year says little about unemployment; its
change says more. This chapter lets the price level move: it builds
\omterm{def:b1:aggregate-supply:sras}{aggregate supply}, joins it to the demand side of the last chapter, and
derives the relation between inflation, expected inflation and
unemployment that central banks still use.

\begin{openingchart}{The change in inflation from the year before against the
unemployment rate, United States, 1960--2018, with the fitted line;
inflation is the percent change of the annual average CPI-U.}{US Bureau of
Labor Statistics, CPI-U all items; unemployment: Jord\`a--Schularick--Taylor
Macrohistory Database, release 6}
\begin{tikzpicture}
\begin{axis}[omscatter, xlabel={Unemployment rate (\%)}, ylabel={Change in inflation (points)}, xmin=2.5, xmax=10.5,
  ymin=-6, ymax=6, width=0.6\linewidth, height=0.45\linewidth]
  \addplot[only marks, mark=*, mark size=1.5pt, omDef] table[x=unemp, y=dpi, col sep=comma] {figdata/bachelor-1/25-aggregate-supply/out/accel.csv};
  \addplot[omThm, thick, sharp plot, mark=none] table[x=u, y=fit, col sep=comma] {figdata/bachelor-1/25-aggregate-supply/out/accel_fit.csv};
  \addplot[only marks, mark=none, nodes near coords, point meta=explicit symbolic, nodes near coords style={font=\scriptsize, anchor=south west, inner sep=1.5pt}]
    table[x=unemp, y=dpi, meta=year, col sep=comma] {figdata/bachelor-1/25-aggregate-supply/out/accel_labels.csv};
\end{axis}
\end{tikzpicture}
\end{openingchart}

The fitted line falls by $0.47$ points of inflation per point of
unemployment and crosses zero at $6.07\%$: below that rate inflation tended
to rise, above it to fall. The correlation is $-0.443$. % figdata: accel_stats
That crossing is an estimate of the natural rate of unemployment.

\begin{recall}
The Phillips curve and its breakdown in the 1970s, stagflation, the
natural rate of unemployment, aggregate demand and supply in words, the
output gap, supply and demand shocks, and the sacrifice ratio were met at
school. IS--LM is \cref{ch:b1:is-lm}; the \omterm{def:b1:unemployment-flows:steady-state-rate}{steady-state unemployment rate} is
\cref{def:b1:unemployment-flows:steady-state-rate}.
\end{recall}

\section{Three short-run supply stories}\label{sec:b1:aggregate-supply:stories}

In the long run output is $\bar Y = F(\bar K, \bar L)$, whatever the price level
(\cref{ch:b1:long-run-model}). In the short run, three stories make output
rise when prices are higher than expected.

\emph{Sticky wages.} Nominal wages are set in advance from the expected
price level; if prices turn out higher, the real wage is lower and firms
hire more. \emph{Sticky prices.} Some firms set prices in advance; when
demand is high, those that cannot adjust sell more. \emph{Imperfect
information.} Producers see their own price before the general price
level, and mistake part of a general rise for a rise in their relative
price.

\begin{definition}[Short-run aggregate supply curve]\label{def:b1:aggregate-supply:sras}
\index{short-run aggregate supply curve}\index{aggregate supply}
The \emph{short-run aggregate supply curve} relates output to the price level
for a given expected price level: $Y = \bar Y + a(P - P^e)$, with $a > 0$. In the long
run $P = P^e$ and the curve is vertical at $\bar Y$.
\end{definition}

\begin{proposition}[Sticky prices give the SRAS]\label{prop:b1:aggregate-supply:sticky}
Suppose a share $1 - s$ of firms set their price flexibly at $p = P + b(Y - \bar Y)$, and a
share $s$ set it in advance at $p = P^e$. Then $P = P^e + \frac{(1 - s)b}{s}(Y - \bar Y)$, that is
$Y = \bar Y + a(P - P^e)$ with $a = s/[(1 - s)b]$.
\end{proposition}

\begin{proof}
The price level is the average, $P = sP^e + (1 - s)[P + b(Y - \bar Y)]$. Subtract $(1 - s)P$:
$sP = sP^e + (1 - s)b(Y - \bar Y)$, and divide by $s$.
\end{proof}

The more firms keep their prices fixed (larger $s$), the flatter the curve:
demand then moves output more and prices less.

\section{AD--AS}\label{sec:b1:aggregate-supply:adas}

\begin{definition}[Aggregate demand curve]\label{def:b1:aggregate-supply:ad}
\index{aggregate demand curve}
The \emph{aggregate demand curve} gives, for each price level, the output at
which the goods and money markets are both in \omterm{def:b1:models-and-data:equilibrium}{equilibrium}: the IS--LM
\omterm{def:b1:models-and-data:equilibrium}{equilibrium} output when $M/P$ varies with $P$.
\end{definition}

A higher $P$ lowers \omterm{def:b1:money-and-inflation:money-demand}{real money balances}, shifts LM up, raises the interest
rate and lowers output: the \omterm{def:b1:aggregate-supply:ad}{aggregate demand curve} slopes down. A rise in
$G$ or $M$ shifts it right.

\begin{omfigure}
\begin{tikzpicture}
\begin{axis}[omrecord, grid=none, xlabel={Output $Y$}, ylabel={Price level $P$}, xmin=80, xmax=120, ymin=0.8,
  ymax=1.25, width=0.56\linewidth, height=0.42\linewidth, samples=60]
  \addplot[demand, domain=82:120] {100/x};
  \addplot[shifted, domain=86:120] {108/x};
  \addplot[supply, domain=80:120] {1 + (x - 100)/200};
  \addplot[supply, dashed, domain=80:120] {1.08 + (x - 100)/200};
  \addplot[black, thick] coordinates {(100,0.8) (100,1.25)};
  \node[font=\footnotesize, anchor=south] at (axis cs:100,1.25) {LRAS};
  \node[font=\footnotesize, anchor=west] at (axis cs:114,1.06) {SRAS};
  \node[font=\footnotesize, anchor=west] at (axis cs:113,0.84) {AD};
  \node[font=\footnotesize, anchor=south west] at (axis cs:113,0.955) {AD$'$};
\end{axis}
\end{tikzpicture}

{\small A demand shock in AD--AS. AD shifts right (dashed blue): in the short
run output and prices rise along SRAS (red). As expected prices catch up,
SRAS shifts up (dashed red) until output is back at $\bar Y = 100$ on the long-run
curve, with prices $8\%$ higher.}
\end{omfigure}

A supply shock --- a rise in the price of oil --- shifts SRAS up: output
falls and prices rise at once, the combination the 1970s called
stagflation. Demand policy can then hold output or hold prices, not both.

\section{The expectations-augmented Phillips curve}\label{sec:b1:aggregate-supply:phillips}

\begin{definition}[Okun's law]\label{def:b1:aggregate-supply:okun}
\index{Okun's law}
\emph{Okun's law} is the empirical rule that output falls below its natural
level by a roughly constant multiple of the excess of unemployment over its
natural rate: $(Y - \bar Y)/\bar Y = -\omega(u - u^n)$.
\end{definition}

In the American data after the Second World War, each point by which
unemployment fell was associated with real output about $3\%$ higher, within
the range of unemployment rates from 3 to 7.5\% (Okun's estimate on data to
1960). % ledger: b1:okun-coefficient

\begin{definition}[Expectations-augmented Phillips curve]\label{def:b1:aggregate-supply:phillips}
\index{Phillips curve, expectations-augmented}\index{expectations-augmented Phillips curve}
The \emph{expectations-augmented Phillips curve} relates inflation to
expected inflation and the unemployment gap: $\pi = \pi^e - \beta(u - u^n) + v$, with $\beta > 0$
and $v$ a supply shock. (Here $\beta$ is the slope of the curve, not a discount
factor.)
\end{definition}

It follows from \omterm{def:b1:aggregate-supply:sras}{aggregate supply}: writing $P - P^e$ in logarithms as $\pi - \pi^e$ (the
price level of last year is known), $\pi = \pi^e + (Y - \bar Y)/a$, and \omterm{def:b1:aggregate-supply:okun}{Okun's law}
replaces $Y - \bar Y$ by $-\omega\bar Y(u - u^n)$, so $\beta = \omega\bar Y/a$ (\cref{exo:b1:aggregate-supply:derive}).
There is a trade-off between inflation and unemployment only for given
expectations. When $u = u^n$ and there are no shocks, $\pi = \pi^e$, at any rate of
inflation: in the long run the curve is vertical.

\section{Adaptive expectations}\label{sec:b1:aggregate-supply:adaptive}

\begin{definition}[Adaptive expectations]\label{def:b1:aggregate-supply:adaptive}
\index{adaptive expectations}
Expectations are \emph{adaptive} when expected inflation is formed from
past inflation; in the simplest form, $\pi^e_t = \pi_{t-1}$.
\end{definition}

\begin{proposition}[Accelerating inflation]\label{prop:b1:aggregate-supply:accelerate}
With $\pi^e_t = \pi_{t-1}$ and no shocks, $\pi_t - \pi_{t-1} = -\beta(u_t - u^n)$. Holding unemployment
below $u^n$ by $d$ points makes inflation rise by $\beta d$ every year, without limit.
\end{proposition}

\begin{proof}
Substitute $\pi^e_t = \pi_{t-1}$ into the Phillips curve; with $u_t = u^n - d$, $\pi_t = \pi_{t-1} + \beta d$, so
$\pi_t = \pi_0 + t\beta d$.
\end{proof}

This is what the opening chart tests: the change in inflation against
unemployment, with the slope $-\beta$ and the crossing at $u^n$. Since the mid-1990s
the slope has been flatter, $-0.134$ over 1996--2018, with a correlation of
only $-0.176$. % figdata: accel_stats
Inflation stopped following its own past.

\begin{model}{The expectations-augmented Phillips curve}
\assumes Expectations formed from past inflation; a natural rate of
unemployment set by real factors; a stable slope.
\predicts Accelerating inflation when $u < u^n$; no long-run trade-off; a
disinflation that costs unemployment, the same total whatever its speed.
\tested The stable trade-off of the 1960s broke in the 1970s; the
change in inflation against unemployment in the United States,
1960--2018, gives a slope of $-0.47$ and a natural rate of $6.07\%$.
% figdata: accel_stats
Across 65 disinflations, the cost per point of inflation was lower when
disinflation was fast and wages flexible (Ball 1994). % ledger: b1:ball1994
\breaks After the mid-1990s inflation stopped responding to its past
(expectations anchored); the curve flattened; and in 2021--23 inflation
moved fast with unemployment low and stable.
\end{model}

\begin{measure}[Estimating the natural rate]
The crossing of the opening chart's line, $6.07\%$, is an estimate over sixty
years. % figdata: accel_stats
Official estimates let the natural rate change over time, as
demographics, matching and benefits change $s$ and $f$
(\cref{ch:b1:unemployment-flows}); they are revised as data arrive. A
natural rate is never observed, only inferred from what inflation did.
\end{measure}

\section{The cost of disinflation}\label{sec:b1:aggregate-supply:disinflation}

With \omterm{def:b1:aggregate-supply:adaptive}{adaptive expectations}, bringing inflation down by $\Delta\pi$ points requires
$\sum_t(u_t - u^n) = \Delta\pi/\beta$ point-years of excess unemployment, whatever the path: each
year's fall in inflation is $\beta$ times that year's excess unemployment. With
Okun's coefficient $\omega$, the output lost is $\omega\Delta\pi/\beta$, and the sacrifice ratio is $\omega/\beta$.

\begin{omfigure}
\begin{tikzpicture}
\begin{axis}[omrecord, grid=none, xlabel={Unemployment rate $u$ (\%)}, ylabel={Inflation $\pi$ (\%)},
  xmin=3, xmax=18, ymin=0, ymax=12, width=0.56\linewidth, height=0.42\linewidth]
  \addplot[demand, domain=3:18] {10 - 0.5*(x - 5)};
  \addplot[demand, dashed, domain=3:18] {4 - 0.5*(x - 5)};
  \addplot[black, thick] coordinates {(5,0) (5,12)};
  \addplot[omThm, thick, mark=*, mark size=1.5pt] coordinates {(5,10) (7,9) (7,8) (7,7) (7,6) (7,5) (7,4) (5,4)};
  \addplot[omThm, thin, dotted, mark=square*, mark size=1.5pt] coordinates {(5,10) (17,4) (5,4)};
  \node[font=\footnotesize, anchor=south] at (axis cs:5,12) {long run};
  \node[font=\footnotesize, anchor=south west] at (axis cs:13.5,6.3) {$\pi^e = 10$};
  \node[font=\footnotesize, anchor=south west] at (axis cs:8.3,2.4) {$\pi^e = 4$};
\end{axis}
\end{tikzpicture}

{\small Two disinflations from $10\%$ to $4\%$ in the weekend problem ($\beta = 0.5$,
$u^n = 5\%$): gradual, one point a year at $u = 7\%$ for six years (solid), or cold turkey,
$u = 17\%$ for one year, along the short-run curve for $\pi^e = 10$, then back to $u^n$
(squares). Both cost 12 point-years of unemployment above
$u^n$. The blue lines are the short-run curves for expected inflation of 10 and
4.}
\end{omfigure}

\begin{record}[The 1970s]
The stable trade-off of the 1960s promised lower unemployment for a little
more inflation. In the 1970s both rose: in 1974 American inflation rose by
$4.82$ points with unemployment at $5.6\%$, and in 1975 unemployment reached
$8.5\%$ with inflation still at $9.13\%$. % figdata: accel_points
With expected inflation rising, the short-run curve had moved up; the
stable trade-off was refuted.
\end{record}

\begin{record}[The Volcker disinflation]
In October 1979 the Federal Reserve changed its operating procedure to
control money growth more firmly, and the federal funds rate reached a
record $20\%$ in late 1980. % ledger: b1:fed-targets
American inflation was $13.5\%$ in 1980 and $3.21\%$ in 1983; unemployment was
$9.7\%$ in 1982 and $9.6\%$ in 1983. % figdata: accel_points
Across 65 disinflations in moderate-inflation countries, fast disinflations
cost less output per point of inflation than gradual ones (Ball 1994).
% ledger: b1:ball1994
The \omterm{def:b1:aggregate-supply:adaptive}{adaptive} model says the speed does not matter; the record says it
does, as if expectations moved faster when the policy was clear.
\end{record}

How much credibility saves depends on how expectations are formed, the
subject of the second university year. If an announcement lowers expected
inflation at once, part of the disinflation is free.

\section{The record}\label{sec:b1:aggregate-supply:record}

\begin{omchart}{Inflation (HICP, annual average) against the unemployment
rate, euro area, 2009--2024, as a path.}{Eurostat (prc\_hicp\_aind,
une\_rt\_a)}
\begin{tikzpicture}
\begin{axis}[omscatter, xlabel={Unemployment rate (\%)}, ylabel={Inflation (\%)}, xmin=6, xmax=13, ymin=-0.5,
  ymax=9.5, width=0.58\linewidth, height=0.42\linewidth]
  \addplot[omDef, thin, sharp plot, mark=*, mark size=1.4pt] table[x=u, y=pi, col sep=comma] {figdata/bachelor-1/25-aggregate-supply/out/euro.csv};
  \addplot[only marks, mark=none, nodes near coords, point meta=explicit symbolic, nodes near coords style={font=\scriptsize, anchor=south west, inner sep=1.5pt}]
    table[x=u, y=pi, meta=year, col sep=comma] {figdata/bachelor-1/25-aggregate-supply/out/euro_points.csv};
\end{axis}
\end{tikzpicture}
\end{omchart}

From 2013 to 2019 euro-area unemployment fell from $12.2\%$ to $7.6\%$ while
inflation stayed between $0.2\%$ and $1.8\%$, with no upward trend: a flat
curve. In 2022 inflation
reached $8.4\%$ with unemployment at $6.7\%$, and in 2024 it was back to $2.4\%$
with unemployment at $6.4\%$. % figdata: euro_points
Such a rise and fall, with almost no movement in unemployment, is
consistent with supply shocks ($v$) and with expectations that stayed
anchored; how much demand added is still debated
(\cref{ch:b1:stabilisation-policy}).

\begin{record}[2021--23]
American consumer prices rose by $9.1\%$ over the twelve months to June 2022,
the peak of the episode, and by $2.9\%$ over the twelve months to December
2024. % figdata: cpi_peak
The \omterm{def:b1:aggregate-supply:adaptive}{adaptive} model would have needed years of high unemployment to bring
inflation down by six points; unemployment barely rose.
\end{record}

\section*{Exercises}\label{sec:b1:aggregate-supply:exercises}

\begin{exercise}[$\star$]\label{exo:b1:aggregate-supply:sras-ex}
With $Y = \bar Y + a(P - P^e)$, $\bar Y = 10{,}000$, $a = 2000$, $P^e = 1$ and $P = 1.05$, find output.
\end{exercise}

\begin{exercise}[$\star$]\label{exo:b1:aggregate-supply:read}
In the opening chart's table, the change in inflation was $4.82$ points in
1974 and $-4.16$ in 1982, with unemployment at $5.6\%$ and $9.7\%$.
% figdata: accel_points
Place each year relative to the fitted line's crossing.
\end{exercise}

\begin{exercise}[$\star$]\label{exo:b1:aggregate-supply:predict}
With $\pi^e = \pi_{-1} = 3\%$, $u = 4\%$, $u^n = 5\%$ and $\beta = 0.5$, predict inflation this year.
\end{exercise}

\begin{exercise}[$\star$]\label{exo:b1:aggregate-supply:okun-ex}
With an Okun coefficient of 2, unemployment is 1 point above the natural
rate. Find the output gap.
\end{exercise}

\begin{exercise}[$\star\star$]\label{exo:b1:aggregate-supply:five}
Starting from $\pi_{-1} = 3\%$, unemployment is held 1 point below $u^n$ for five years, with
$\beta = 0.5$ and \omterm{def:b1:aggregate-supply:adaptive}{adaptive expectations}. Give the inflation path.
\end{exercise}

\begin{exercise}[$\star\star$]\label{exo:b1:aggregate-supply:sacrifice}
Inflation is cut from $10\%$ to $4\%$ with $\beta = 0.5$ and an Okun coefficient of 2.
Compute the point-years of excess unemployment, the output lost and the
sacrifice ratio.
\end{exercise}

\begin{exercise}[$\star\star$]\label{exo:b1:aggregate-supply:estimate}
The fitted line of the opening chart is $\Delta\pi = 2.853 - 0.47u$. % figdata: accel_stats
Estimate $u^n$ and $\beta$.
\end{exercise}

\begin{exercise}[$\star\star$]\label{exo:b1:aggregate-supply:supply}
A supply shock $v = 2$ hits an economy at $u = u^n$ with $\pi^e = \pi_{-1}$ and $\beta = 0.5$.
Compare holding unemployment at $u^n$ with holding inflation at its previous
level.
\end{exercise}

\begin{exercise}[$\star\star$]\label{exo:b1:aggregate-supply:demand}
In AD--AS, the money supply rises by $8\%$. Describe the short-run and
long-run effects on output and prices.
\end{exercise}

\begin{exercise}[$\star\star\star$]\label{exo:b1:aggregate-supply:derive}
Derive the \omterm{def:b1:aggregate-supply:phillips}{expectations-augmented Phillips curve} from the sticky-price SRAS
and \omterm{def:b1:aggregate-supply:okun}{Okun's law}, and express $\beta$ in terms of $a$ and $\omega$.
\end{exercise}

\begin{exercise}[$\star\star\star$]\label{exo:b1:aggregate-supply:lambda}
Let $\pi^e_t = \lambda\pi_{t-1} + (1 - \lambda)\bar\pi$, with $0 \le \lambda \le 1$. Show that a permanent trade-off between
inflation and unemployment exists unless $\lambda = 1$, and say what the 1970s
suggest about $\lambda$.
\end{exercise}

\begin{exercise}[$\star\star\star$]\label{exo:b1:aggregate-supply:refute}
What pattern after 1995 would refute \omterm{def:b1:aggregate-supply:adaptive}{adaptive expectations}? Does the slope
of $-0.134$ over 1996--2018 show it? % figdata: accel_stats
\end{exercise}

\section{Problem: Kestria's Disinflation}\label{sec:b1:aggregate-supply:problem}

\begin{problem}[{Weekend problem --- the cost of bringing inflation down}]\label{pb:b1:aggregate-supply:kestria}
Kestria (invented) has inflation of $10\%$ and wants $4\%$. Its Phillips curve
is $\pi_t = \pi^e_t - 0.5(u_t - 5)$, with expectations $\pi^e_t = \pi_{t-1}$ unless stated, and
\omterm{def:b1:aggregate-supply:okun}{Okun's law} gives an output gap of $-2(u - 5)\%$.

\medskip\noindent\textbf{Part I --- The curve.}
\begin{enumerate}
  \item What is the natural rate?
  \item What unemployment keeps inflation at $10\%$?
  \item If unemployment were held at $3\%$, give inflation for five years.
  \item What is the sacrifice ratio implied by the two coefficients?
  \item Why is there no long-run trade-off?
\end{enumerate}

\medskip\noindent\textbf{Part II --- Cold turkey.} Inflation is brought to $4\%$ in one
year and held there.
\begin{enumerate}[resume]
  \item Find unemployment in the first year.
  \item Find unemployment in the following years.
  \item Compute the point-years of excess unemployment.
  \item Compute the output lost.
  \item What makes this path politically hard?
\end{enumerate}

\medskip\noindent\textbf{Part III --- Gradualism.}
\begin{enumerate}[resume]
  \item Inflation falls one point a year. Find unemployment each year.
  \item Compute the point-years of excess unemployment.
  \item In two steps, to $7\%$ then $4\%$: find unemployment each year.
  \item Show that every path costs the same under \omterm{def:b1:aggregate-supply:adaptive}{adaptive expectations}.
  \item What does Ball's evidence on 65 disinflations say about that result?
  \item Why might the cost depend on the speed in reality?
\end{enumerate}

\medskip\noindent\textbf{Part IV --- Credibility.} The central bank announces the
target, and the announcement lowers expected inflation in the first year
by 4 points.
\begin{enumerate}[resume]
  \item Find expected inflation in the first year.
  \item Find the unemployment needed to reach $4\%$ at once.
  \item Compute the point-years of excess unemployment.
  \item Compute the saving.
  \item Compute the output lost and compare.
  \item What would make the announcement believed?
  \item What happens if the bank announces but does not act?
  \item How does the 2021--23 episode bear on credibility?
  \item State the result: the cost of the cold-turkey and gradual paths, and what credibility saves.
\end{enumerate}
\end{problem}
